% Phase 6D — readable draft with verified bibliography (M33 TDF tau-geometry)
% Numerical values: outputs/tables/phase6a_key_results_table.csv (fit mask 0.4--23 kpc).
% Bibliography: Corbelli2014 and LopezFune2017 verified — see docs/manuscript_bibliography_verification.md

\documentclass[11pt,a4paper]{article}
\usepackage[margin=2.5cm]{geometry}
\usepackage{graphicx}
\usepackage{booktabs}
\usepackage{amsmath}
\usepackage{hyperref}
\usepackage{array}
\hypersetup{colorlinks=true,linkcolor=blue,citecolor=blue}

% Figures optional until pipeline PNGs are copied beside the draft.
\newcommand{\includefig}[2][]{%
  \IfFileExists{#2}{%
    \includegraphics[#1]{#2}%
  }{%
    \fbox{\parbox[c][0.2\linewidth][c]{0.92\linewidth}{\centering\small
      [Figure placeholder]\\[0.4em]
      \texttt{\detokenize{#2}}}}}%
}

\title{Reconstructing the Missing Acceleration of M33 as a Low-Parameter Time-Delay Field Geometry}
\author{M33 TDF $\tau$-Geometry Collaboration\\ \textit{Phase 6D draft --- claim review required; bibliography verified for cited sources}}
\date{\today}

\begin{document}
\maketitle

\begin{abstract}
We present a reproducible analysis pipeline for the nearby spiral galaxy M33 as an \textbf{intermediate-scale test} of whether the missing acceleration inferred from rotation dynamics can be reconstructed as a smooth, low-parameter time-delay field ($\tau$) geometry, and whether the \textbf{same reconstructed $\tau$-map} yields \textbf{normalized deflection-proxy} predictions without adding a separate halo component or performing a lensing-only $\tau$ fit.
Using the Corbelli et al.\ (2014) rotation curve and an in-repository baryonic decomposition (quality flag PASS\_WITH\_CAVEAT), we compare a baryonic-only diagnostic, standard NFW and Burkert halo baselines, and low-parameter $\tau$-gradient knot models on a fixed fit mask ($0.4$--$23$\,kpc, $n{=}56$).
The low-parameter three-knot TDF model is \textbf{competitive with the NFW baseline} on this processed M33 dataset (RMSE $\approx 3.51$\,km\,s$^{-1}$, AIC $\approx 63.0$ vs.\ NFW AIC $\approx 68.0$, BIC $\approx 72.0$).
A five-knot variant attains the lowest RMSE ($\approx 3.35$\,km\,s$^{-1}$) but incurs a higher BIC penalty ($\approx 76.5$).
Phase~3D sensitivity audits yield a PASS\_WITH\_CAVEAT stability conclusion under tested mask and $K_\tau$ variants.
An axisymmetric two-dimensional $\tau$-map and Corbelli-based sky projection are generated from the frozen radial solution; Phase~5A produces a \texttt{normalized\_proxy} deflection field from that map.
A L\'opez Fune et al.\ (2017) comparison of an enclosed-mass-style proxy $M_{\tau,\mathrm{eff}}=r v_\tau^2/G$ to literature $M_{\mathrm{enclosed}}$ at $\sim 23$\,kpc gives PASS\_WITH\_CAVEAT (ratio $\approx 1.12$) as a \textbf{dynamical upper-bound consistency check}, not weak lensing.
We do \emph{not} claim that dark matter is disproven, that TDF replaces $\Lambda$CDM, or that lensing is confirmed; physical arcsecond calibration and M33 weak-lensing data remain future work.
\end{abstract}

\section{Introduction}
\label{sec:intro}

Rotation curves of disk galaxies remain a central diagnostic of mass distribution.
For M33---a gas-rich Local Group spiral---the observed rotation curve exceeds what visible baryons typically explain when using standard disk-halo decompositions \citep{Corbelli2014}.
In $\Lambda$CDM interpretations, this ``missing'' acceleration is attributed to a dark-matter halo, often parameterized with NFW or cored profiles.
Alternative organizing pictures ask whether the missing acceleration can be described by additional dynamical structure in the baryonic problem itself, rather than introducing a separate collisionless component in a given analysis branch.

The time-delay field (TDF) framework treated here introduces a radial $\tau$-geometry coupled to baryons through a low-parameter gradient model.
The scientific question is intentionally narrow: on a \textbf{single galaxy} and a \textbf{documented processed dataset}, can we (i) represent rotation residuals with a smooth, low-parameter $\tau$ model competitive with a standard halo baseline, and (ii) propagate the \textbf{same frozen $\tau$-map} to a deflection \emph{proxy} without refitting $\tau$ for lensing or adding a separate halo?

This manuscript reports results from a version-controlled pipeline (Phases~1D--6C).
We emphasize \textbf{consistency and claim control} over cosmological inference.
The analysis does not constitute evidence that dark matter is disproven, does not replace $\Lambda$CDM, and does not confirm gravitational lensing in physical units.

\section{Data and provenance}
\label{sec:data}

\subsection{Rotation sample}

The canonical processed table is \texttt{data/processed/m33\_rotation.csv}, containing 58 radii with observed velocities from Corbelli et al.\ (2014) Table~1 \citep{Corbelli2014}.
Each row records $r$, $v_{\mathrm{obs}}$, uncertainty, and derived baryonic components used downstream.
Provenance, extraction steps, and validation are documented in \texttt{docs/data\_sources.md} and enforced by \texttt{scripts/validate\_m33\_data.py}.

\subsection{Fit mask and comparison design}

Fair comparison among baryonic-only, halo, and TDF models uses a common fit mask $0.4 \le r \le 23$\,kpc, yielding $n{=}56$ points.
This mask matches the Phase~2C/3C audit convention and the outer radius of the low-parameter TDF fits.
Metrics quoted in Table~\ref{tab:models} are taken from \texttt{outputs/tables/phase6a\_key\_results\_table.csv}, which consolidates \texttt{phase3c\_combined\_model\_comparison.csv}.

\subsection{Geometry for projection}

Sky-plane projection (Phase~4B) uses Corbelli et al.\ (2014) tilted-ring geometry (\texttt{corbelli2014\_tilted\_ring\_geometry\_model\_shape.csv}), not ad hoc placeholder inclinations.

\begin{table}[ht]
  \centering
  \caption{Data provenance summary (Table~1). Repository paths; expand before journal submission.}
  \label{tab:provenance}
  \begin{tabular}{>{\raggedright\arraybackslash}p{0.22\linewidth}p{0.34\linewidth}p{0.34\linewidth}}
    \toprule
    Data product & Literature / method & Repository artifact \\
    \midrule
    Observed rotation & Corbelli et al.\ 2014 Table~1 & \texttt{m33\_rotation.csv} \\
    Baryonic $v_{\mathrm{gas}}$, $v_{\mathrm{disk}}$ & Phase~1D-D1 disk gravity & same; PASS\_WITH\_CAVEAT \\
    Tilted-ring geometry & Corbelli Fig.~3 / Appendix & extracted CSV \\
    Model metrics & Phases~2C, 3C, 6A & \texttt{phase6a\_key\_results\_table.csv} \\
    \bottomrule
  \end{tabular}
\end{table}

\section{Baryonic decomposition and caveats}
\label{sec:baryons}

Baryonic velocities are \textbf{not} taken directly from published Corbelli decomposition columns.
Instead, gas and stellar disk contributions are derived from Tabulated surface densities using the documented axisymmetric pipeline (Phase~1D-D1), including helium and molecular gas prescriptions stated in \texttt{docs/data\_sources.md}.

A corrected Fig.~12 sanity comparison yields PASS\_WITH\_CAVEAT: derived velocities can differ from spot-checked literature curves by approximately $7$--$10$\,km\,s$^{-1}$ at some radii.
Digitized Fig.~12 curves are \textbf{not} used as canonical velocities.

\textbf{Implication.} Every halo and TDF fit in this paper inherits the baryonic caveat.
Comparative statements (e.g.\ TDF vs.\ NFW) are conditional on this processed baryonic model, not on an independent third-party decomposition.

\section{Baseline models}
\label{sec:baselines}

We report three baselines before the TDF branch: baryonic-only, NFW, and Burkert.
All use the same mask and D1 baryonic inputs.

\subsection{Baryonic-only diagnostic}

With no halo or $\tau$ term, the baryonic-only curve fails to match observations on the fit mask: RMSE $\approx 58.6$\,km\,s$^{-1}$, $\chi^2 \approx 2.0\times 10^4$, reduced $\chi^2 \approx 365$ (Table~\ref{tab:models}).
This supports the standard statement that \textbf{baryons alone do not explain the rotation curve} on this processed dataset, while noting the derived-baryonic caveat above.

\subsection{NFW ($\Lambda$CDM comparison baseline)}

A two-parameter NFW halo fit yields RMSE $\approx 4.16$\,km\,s$^{-1}$, reduced $\chi^2 \approx 1.21$, AIC $\approx 68.0$, and BIC $\approx 72.0$.
We treat NFW as a strong \textbf{$\Lambda$CDM comparison baseline} on D1 baryons---not as validation of TDF microphysics and not as a refutation of alternative frameworks.

\subsection{Burkert (boundary-limited here)}

The Burkert fit achieves RMSE $\approx 23.0$\,km\,s$^{-1}$ with AIC $\approx 2342$.
In the Phase~2C audit, the core radius often sits at the \textbf{upper fit bound} ($r_0 \approx 200$\,kpc), with mismatches versus Corbelli BVI reference scales documented in \texttt{phase2c\_model\_audit\_summary.csv}.
We therefore describe Burkert as \textbf{boundary-limited in this processed dataset} and avoid using it as a preferred publication anchor.

\begin{figure}[ht]
  \centering
  \includefig[width=0.88\linewidth]{../outputs/figures/phase2a_baryonic_only_rotation_curve.png}
  \caption{Figure~1: Observed M33 rotation curve compared to the baryonic-only model (Phase~2A). Command: \texttt{run\_phase2a\_baryonic\_only.py}.}
  \label{fig:rot-baryon}
\end{figure}

\section{TDF radial reconstruction}
\label{sec:tdf-radial}

The TDF branch begins from the squared residual acceleration proxy $\Delta v^2 = v_{\mathrm{obs}}^2 - v_{\mathrm{bar}}^2$ on the processed table.
Phases~3A and~3B implement direct and regularized reconstructions of $\tau$-gradients with fixed $K_\tau=1$.

\textbf{Important distinction.} Phase~3A/3B reconstructions are useful diagnostics of $\Delta v^2$ structure and smoothing trade-offs; they are \textbf{not} fair low-parameter competitors in AIC/BIC.
The formal model comparison uses Phase~3C knot models with a small, explicit parameter count.

The velocity-squared contribution follows the project convention $v_\tau^2 = r\,K_\tau\,\mathrm{d}\tau/\mathrm{d}r$ with $K_\tau=1$ as a \textbf{normalization convention}, not an independently calibrated physical constant (see Phase~3D).

\section{Low-parameter TDF comparison}
\label{sec:tdf-compare}

\subsection{Model family}

We fit piecewise-linear $\tau$-gradient knot values for $k\in\{3,4,5\}$ with $K_\tau$ fixed (\texttt{tdf\_lowparam\_3knot}, etc.).
The primary science model for maps and lensing is the \textbf{three-knot} solution, selected for balance of fit quality and parameter economy.

\subsection{Metrics}

On the fit mask, the three-knot model attains RMSE $\approx 3.51$\,km\,s$^{-1}$, $\chi^2 \approx 57.0$, AIC $\approx 63.0$, and BIC $\approx 69.0$.
The five-knot model achieves the lowest RMSE ($\approx 3.35$\,km\,s$^{-1}$, $\chi^2 \approx 56.3$) but higher AIC/BIC ($\approx 66.3$ / $76.5$) due to parameter penalty.

\textbf{Conservative statement.} The low-parameter three-knot TDF model is \textbf{competitive with the NFW baseline} on this processed M33 dataset under AIC/BIC on the standard mask.
This is a rotation-dynamics consistency statement only.

\begin{figure}[ht]
  \centering
  \includefig[width=0.88\linewidth]{../outputs/figures/phase3c_tdf_lowparam_rotation_comparison.png}
  \caption{Figure~2: Rotation-curve comparison for NFW, Burkert, and TDF models (Phase~3C).}
  \label{fig:rot-compare}
\end{figure}

\begin{figure}[ht]
  \centering
  \includefig[width=0.88\linewidth]{../outputs/figures/phase3c_tdf_lowparam_tau_gradient.png}
  \caption{Figure~3: Reconstructed $\tau$-gradient and $\tau$-profile for \texttt{tdf\_lowparam\_3knot}.}
  \label{fig:tau-profile}
\end{figure}

\begin{table}[ht]
  \centering
  \caption{Model comparison on fit mask $0.4$--$23$\,kpc (Table~2; \texttt{phase6a\_key\_results\_table.csv}).}
  \label{tab:models}
  \begin{tabular}{lrrrrr}
    \toprule
    Model & $n_{\mathrm{fit}}$ & RMSE [km\,s$^{-1}$] & $\chi^2$ & AIC & BIC \\
    \midrule
    Baryonic only & 56 & 58.56 & 20096 & 20096 & 20096 \\
    NFW & 56 & 4.16 & 64.0 & 68.0 & 72.0 \\
    Burkert & 56 & 23.02 & 2338 & 2342 & 2346 \\
    TDF 3-knot & 56 & 3.51 & 57.0 & 63.0 & 69.0 \\
    TDF 5-knot & 56 & 3.35 & 56.3 & 66.3 & 76.5 \\
    \bottomrule
  \end{tabular}
\end{table}

\section{Sensitivity and stability}
\label{sec:sensitivity}

Phase~3D (\texttt{outputs/tables/phase3d\_tdf\_sensitivity\_summary.csv}) audits knot count, $K_\tau$ sweeps, Gaussian smoothing (diagnostic), and fit-mask variants.

Findings relevant to manuscript language:
\begin{itemize}
  \item \texttt{tdf\_lowparam\_3knot} is the best TDF model by AIC/BIC among knot counts, beating NFW on AIC at the default mask.
  \item $K_\tau$ exhibits expected \textbf{scale degeneracy} when knots are refit; velocity RMSE spread across the sweep is modest ($\max \approx 0.42$\,km\,s$^{-1}$) but $K_\tau$ must not be over-interpreted as a measured constant.
  \item Stricter masks (e.g.\ $1.0$--$22.0$\,kpc) can change AIC rankings; we summarize this as \textbf{PASS\_WITH\_CAVEAT} stability, not universal robustness.
\end{itemize}

We do not claim that TDF is insensitive to all systematics or independent of baryonic derivation choices.

\section{Two-dimensional $\tau$-map and sky projection}
\label{sec:2dmap}

\subsection{Disk-plane map (Phase~4A)}

The two-dimensional map is an \textbf{axisymmetric extension} of the frozen three-knot radial profile:
$\tau_{2\mathrm{D}}(x,y)=\tau_{\mathrm{radial}}(R)$.
No additional spatial parameters are fitted; spiral arms and non-axisymmetric gas features are not modeled.

\subsection{Sky projection (Phase~4B)}

The disk-plane map is projected to sky coordinates using Corbelli tilted-ring geometry, producing \texttt{outputs/maps/phase4b\_tau\_sky\_projected\_map.npz}.
Metadata in \texttt{phase4a\_tau\_2d\_map\_metadata.csv} and \texttt{phase4b\_tau\_projection\_metadata.csv} records source model \texttt{tdf\_lowparam\_3knot} and geometry flags.

\begin{figure}[ht]
  \centering
  \includefig[width=0.44\linewidth]{../outputs/figures/phase4a_tau_2d_map.png}
  \hfill
  \includefig[width=0.44\linewidth]{../outputs/figures/phase4b_tau_sky_projected_map.png}
  \caption{Figure~4: Disk-plane $\tau$ map (left) and sky-projected $\tau$ (right).}
  \label{fig:2dmap}
\end{figure}

\section{Deflection proxy and upper-bound consistency}
\label{sec:lensing}

\subsection{Normalized deflection proxy (Phase~5A)}

From the sky-plane $\tau$-map, Phase~5A constructs deflection and convergence \textbf{proxies} in \texttt{normalized\_proxy} units ($\alpha_\tau$ scale $=1$).
\textbf{The same reconstructed $\tau$-map produces normalized deflection-proxy fields without adding a separate halo.}
$\tau$ is not refit for lensing; no lensing-only degrees of freedom are introduced.

This is a \textbf{prediction scaffold} for future calibration, not a claim of observational lensing detection.
The deflection proxy is \textbf{not} a physical arcsecond prediction.

\begin{figure}[ht]
  \centering
  \includefig[width=0.78\linewidth]{../outputs/figures/phase5a_deflection_magnitude_map.png}
  \caption{Figure~5: Magnitude of the normalized deflection proxy (Phase~5A).}
  \label{fig:deflection}
\end{figure}

\subsection{L\'opez Fune dynamical upper bound (Phase~5C-B)}

We compare a Newtonian-equivalent enclosed-mass proxy $M_{\tau,\mathrm{eff}}(<r)=r\,v_\tau^2/G$ (with $G=4.30091\times 10^{-6}$ in kpc\,km$^2$\,M$_\odot^{-1}$) to the L\'opez Fune et al.\ (2017) Burkert enclosed-mass reference $M_{\mathrm{enclosed}}\approx 6.7\pm 1.2\times 10^{10}\,M_\odot$ at $\sim 23$\,kpc \citep{LopezFune2017}.

At the maximum profile radius $r=22.72$\,kpc, the pipeline gives $M_{\tau,\mathrm{eff}}\approx 7.5\times 10^{10}\,M_\odot$ and ratio $M_{\tau,\mathrm{eff}}/M_{\mathrm{L\'opez}}\approx 1.12$.
The audit status is \textbf{PASS\_WITH\_CAVEAT} because $M_{\tau,\mathrm{eff}}$ lies below the $1\sigma$ upper envelope ($\approx 7.9\times 10^{10}\,M_\odot$).

\textbf{The L\'opez Fune comparison is a dynamical upper-bound consistency check.}
It uses rotation-based literature with partial circularity through shared Corbelli inputs; it is \textbf{not} a weak-lensing test of the deflection map.

\section{Limitations}
\label{sec:limits}

\begin{enumerate}
  \item \textbf{Single galaxy.} Results are specific to M33 and this processed table; they do not by themselves establish a general replacement for dark matter or $\Lambda$CDM.
  \item \textbf{Baryonic systematics.} D1-derived velocities carry PASS\_WITH\_CAVEAT; all comparative metrics shift if baryonic inputs change.
  \item \textbf{$K_\tau$ normalization.} $K_\tau=1$ is a convention; Phase~3D documents degeneracy with refit knot amplitudes.
  \item \textbf{Burkert baseline.} Boundary-limited here; not a reliable secondary anchor for M33 in this setup.
  \item \textbf{Axisymmetric 2D map.} No independent 2D $\tau$ fit; not a morphological model of gas or arms.
  \item \textbf{Lensing units.} Deflection remains \texttt{normalized\_proxy}; no arcsecond calibration or source-plane inference.
  \item \textbf{Weak lensing.} No M33-specific weak-lensing map is integrated (\texttt{m33\_direct\_weak\_lensing\_gap}); weak-lensing comparison is future work.
  \item \textbf{Dark matter.} Dark matter is \textbf{not disproven}; NFW remains the reference $\Lambda$CDM baseline for comparison.
\end{enumerate}

\section{Future work}
\label{sec:future}

Near-term extensions, without changing the numerical results reported here:
\begin{itemize}
  \item Document and implement physical $\alpha_\tau$ calibration to arcseconds if independent constraints become available.
  \item Incorporate M33 weak-lensing or other lensing data when a documented map exists in the repository.
  \item Improve baryonic velocities (reduce PASS\_WITH\_CAVEAT severity) and propagate uncertainties through $\tau$ and deflection proxies.
  \item Extend to additional galaxies with the same claim-control discipline.
\end{itemize}

\section{Conclusion}
\label{sec:conclusion}

We have implemented an end-to-end, reproducible M33 pipeline that tests a low-parameter $\tau$-geometry interpretation of missing rotation acceleration and propagates the same frozen map to normalized deflection proxies.

On the standard $0.4$--$23$\,kpc mask, the three-knot TDF model is \textbf{competitive with the NFW baseline} (RMSE $\approx 3.51$\,km\,s$^{-1}$, AIC $\approx 63.0$ vs.\ NFW $\approx 68.0$), while baryons alone fail dramatically (RMSE $\approx 58.6$\,km\,s$^{-1}$).
Sensitivity audits support a \textbf{PASS\_WITH\_CAVEAT} stability summary, not a claim of universal insensitivity.
The same $\tau$-map yields sky-projected fields and a \textbf{normalized deflection-proxy} without a separate halo.
L\'opez Fune enclosed mass provides a \textbf{dynamical upper-bound consistency check} (PASS\_WITH\_CAVEAT, ratio $\approx 1.12$), not lensing confirmation.

We do \emph{not} claim that M33 proves TDF, that lensing is confirmed, or that dark matter is disproven.
Physical lensing calibration and weak-lensing comparison remain necessary before observational lensing statements.

\begin{table}[ht]
  \centering
  \caption{Claim and caveat summary (Table~3; excerpt from \texttt{phase6a\_claim\_traceability\_matrix.csv}).}
  \label{tab:claims}
  \begin{tabular}{>{\raggedright\arraybackslash}p{0.38\linewidth}p{0.14\linewidth}p{0.38\linewidth}}
    \toprule
    Claim & Status & Caveat \\
    \midrule
    3-knot TDF competitive with NFW & supported & Rotation-only; D1 baryons \\
    Same $\tau$-map $\to$ deflection proxy & supported & \texttt{normalized\_proxy}; no separate halo \\
    L\'opez Fune upper-bound check & supported & Dynamics; PASS\_WITH\_CAVEAT; not WL \\
    Weak lensing comparison & future work & No M33 WL map in repo \\
    Dark matter disproven & \textbf{prohibited} & Do not state \\
    \bottomrule
  \end{tabular}
\end{table}

\section*{Reproducibility}
Scripts and phase order are listed in \texttt{outputs/reports/phase6a\_reproducibility\_commands.md}.
Phase~6A--6D audits: \texttt{run\_phase6a\_publication\_audit.py} through \texttt{run\_phase6d\_manuscript\_polish\_audit.py}.

\bibliographystyle{plain}
\begin{thebibliography}{99}
\bibitem[Corbelli et al.(2014)]{Corbelli2014}
Corbelli, E., Thiliker, D.~A., Zschaechner, L., et al.\ 2014,
\textit{Dynamical signatures of a $\Lambda$CDM-halo and the distribution of the baryons in M33},
A\&A, 572, A23.
\href{https://doi.org/10.1051/0004-6361/201424033}{doi:10.1051/0004-6361/201424033}.
\bibitem[L\'opez Fune et al.(2017)]{LopezFune2017}
L\'opez Fune, M.~C., Salucci, P., \& Corbelli, E.\ 2017,
\textit{The radial dependence of dark matter distribution in M33},
MNRAS, 474, 4010.
\href{https://doi.org/10.1093/mnras/stx2742}{doi:10.1093/mnras/stx2742}.
\end{thebibliography}

\end{document}
